\documentclass[../main.tex]{subfiles}
\begin{document}

\section{Conclusion}
\label{sec:conclusion}

This thesis asks when outlyingness translates into joint influence on a prespecified regression coefficient, and how fixed-\(k\) Most Influential Sets characterise the cases in which the two diverge. The results show that the relationship depends on the mechanism generating unusual observations or model departures. Observations can be unusual because of their outcomes, covariates, or relationship to the fitted model while exerting little joint influence on the coefficient of interest. Fixed-\(k\) MIS characterises the complementary quantity: for a prespecified deletion budget, it identifies the set that maximises movement in the target coefficient under the stated fixed-residualised objective. Its interpretation is therefore coefficient-specific and conditional on \(k\).

The distributional-contamination simulation provides the clearest evidence for this distinction. Bad leverage produces strong alignment between the injected contamination set and target-coefficient influence because extreme predictor positions are combined with response shifts that distort the fitted slope. MIS consequently recovers most of the injected set and generally attains higher recovery than the classical observation-level diagnostics. Response outliers and good leverage produce different patterns. Large residuals can be unusual without strongly determining the target slope, while extreme predictor positions can remain closely aligned with the fitted relationship. Lower injected-set recovery by MIS in these designs therefore reflects divergence between contamination status and target-coefficient influence. Recovery of a planted set is informative when the planted mechanism itself generates concentrated influence on the target coefficient.

The robust-estimation simulation establishes a second distinction. Localising coefficient-influential observations and obtaining a stable coefficient estimate under contamination are different statistical objectives. MM regression and LTS modify estimation directly to reduce the effect of contamination, whereas fixed-\(k\) MIS identifies observations according to their joint effect on a prespecified coefficient. Post-deletion OLS can therefore move an otherwise stable coefficient when the deletion budget requires selection of the set producing the largest coefficient change. The coefficient-error distributions also show that typical Monte Carlo behaviour can differ from average performance when infrequent large errors affect arithmetic means more strongly than medians. These results support interpreting MIS as a localisation diagnostic and robust regression as a complementary estimation strategy.

The misspecification simulation extends the analysis beyond injected contamination. Several forms of misspecification generate substantial target-coefficient sensitivity even though no conventional contamination set is introduced, and the response is strongly mechanism dependent. Nonlinear confounding, heterogeneous slopes, nonlinear functional form, and threshold behaviour generate strong MIS responses across the recorded environments, whereas omitted-variable misspecification is substantially weaker in the simulated design. Linear endogeneity provides an exact boundary case. Because its severity component lies within the fitted linear span, increasing severity changes the coefficient representation without changing the standardised deletion-sensitivity structure. The presence of misspecification therefore does not by itself determine whether influence becomes concentrated in a small set of observations.

The same study separates coefficient sensitivity from localisation of a substantively defined affected subgroup. In the threshold design, MIS-selected observations are substantially enriched for membership in the affected group. Heterogeneous slopes produce much weaker enrichment despite substantial target-coefficient sensitivity. A mechanism can therefore produce concentrated coefficient influence without causing the maximising set to correspond closely to a known subgroup. Post-deletion estimation provides a further boundary. Removing selected observations does not restore omitted structure, exogeneity, functional form, or other maintained assumptions, and coefficient error can increase as the deletion budget grows. MIS characterises where finite-sample coefficient sensitivity is concentrated; correction of the underlying model requires a separate analysis.

The empirical applications show that these distinctions also arise in realised samples. Ten published experimental applications are audited under a common coefficient-targeted deletion protocol with full-refit validation. Eight of the ten selected coefficients cross zero along at least one searched deletion path within an approximately \(5\%\) deletion budget. The first zero crossing ranges from one observation, approximately \(0.009\%\) of the estimation sample in the cash-versus-in-kind application, to approximately \(3.46\%\) of the sample in the electricity-information application. The selected sets also differ in their structure. Some contain a persistent influential core as \(k\) increases, others change membership across deletion budgets, and the immunization application exhibits geographic concentration. These patterns show substantial heterogeneity in the dependence of audited coefficients on particular observations in the realised sample.

A central limitation of the empirical analysis concerns interpretation of these influential sets. MIS identifies which observations jointly produce large movement in the audited coefficient, but the deletion result alone does not explain why those observations are influential. Their influence may be associated with unusual outcomes or covariates, heterogeneous effects, geographic or institutional structure, measurement features, features of the sampling process, or model misspecification. Distinguishing among these possibilities requires study-specific information beyond the coefficient-deletion diagnostic. The geographic analysis of the immunization application illustrates one possible follow-up by showing that the first sign-reversing set is unevenly distributed across districts. This spatial concentration is descriptive evidence about the composition of the influential set; it does not establish the mechanism that generates its influence. The empirical applications therefore locate finite-sample coefficient dependence without providing a substantive explanation for its origin.

The empirical findings also remain conditional on the audited estimand and specification. A zero crossing or large coefficient movement establishes sensitivity of that point estimate under the stated deletion budget. It does not determine the validity of the source study's identification strategy, the sensitivity of other estimands, or the appropriate substantive interpretation of the original study. Several applications use documented linear companion or robustness specifications because their headline estimators lie outside the current MIS search framework. Full-refit validation evaluates the selected set using the corresponding frozen audit specification, while the search remains tied to the fixed-residualised linear objective.

\subsection{Limitations and Future Development}

Several additional limitations define the methodological scope of the analysis. The deletion size \(k\) is prespecified and is interpreted as a sensitivity budget rather than an estimated number of anomalous observations. Different values of \(k\) can produce different influential sets, and the method does not select a uniquely correct deletion size. MIS is also coefficient-specific, so different target coefficients can generate different influential sets within the same fitted model. Global optimisation in this thesis applies to the stated fixed-\(k\), fixed-residualised objective; full-refit validation evaluates the selected candidate set but does not exhaustively optimise the refitted coefficient over every possible size-\(k\) deletion. The contamination simulations additionally use oracle-assisted deletion budgets and, in some designs, the known direction of the induced coefficient perturbation. EVT calibration is conditional on the chosen deletion budget and becomes restrictive when block size approaches the deletion-set size.

The empirical limitation suggests a natural direction for future development. Once MIS has located an influential set, subsequent analysis could characterise its members using observable features that are substantively meaningful for the application. These may include baseline characteristics, treatment status, outcomes, residuals, geographic location, institutional setting, or other study-specific variables. Persistence across deletion budgets could also distinguish observations that repeatedly form an influential core from those that enter only at particular values of \(k\). Such analysis would provide evidence about the observable structures associated with coefficient sensitivity while preserving the distinction between localisation and explanation.

A second direction is methodological. Recent work has extended influential-set analysis from adaptive search procedures to global optimisation of fixed-\(k\) linear-regression objectives \parencite{kuschnig2021hidden,konrad2026finding}. Further work could extend coefficient-targeted subset search to structured deletion constraints and to estimands and estimators outside the current linear representation. The empirical applications in this thesis make the latter limitation particularly relevant: several headline specifications use instrumental-variable, GMM, or nonlinear estimators and can currently be audited only through documented linear companion specifications. Extending influential-set methods to these estimators would allow the same coefficient-sensitivity question to be applied directly to a broader class of empirical models while preserving the distinction between influence diagnosis and substantive interpretation of the selected observations.

Fixed-\(k\) MIS is therefore most useful as a targeted diagnostic of joint coefficient influence. The simulations show that outlyingness and target-coefficient influence coincide under some mechanisms and diverge under others, while the empirical applications show that this distinction can be consequential in realised samples. MIS identifies which small set of observations jointly determines a prespecified coefficient under a stated deletion budget. The result does not assign those observations a general outlier status or explain the substantive process that made them influential. Unusual observations need not be coefficient influential, coefficient-influential observations need not correspond to a known anomalous group, and localisation of influence does not itself provide robust estimation, model correction, or causal validation. These distinctions define both the contribution and the appropriate scope of fixed-\(k\) coefficient-targeted influence analysis.

\end{document}
